\section{Notation for the order interpolation}
\label{sec:order-notation}

Write $H_0=0$, $H_n=\sum_{j=1}^n1/j$ and
\begin{equation}\label{sorder:def:S}
 S(s)=\sum_{n=0}^{\infty}\frac{(-1)^nH_n}{(2n+1)^s},\qquad \Re s>0.
\end{equation}
Use the positive real logarithm in the denominator, and let
$\beta(s)=\sum_{n\ge0}(-1)^n(2n+1)^{-s}$ denote the Dirichlet beta
function and its entire continuation. Put
\[
 c=\frac\pi2,\qquad G=\beta(2),\qquad
 L(x)=\gamma+\Re\psi\left(\frac{1+ix}{2}\right)\quad(x\ge0).
\]
The gamma logarithmic derivative is $\psi$, and $\gamma=-\psi(1)$.
In particular, $L(0)=-2\log2$.

Analytic continuation of alternating Euler sums and their values at
negative integers belong to an established subject; see
Boyadzhiev--Gadiyar--Padma \cite{BGP}. Our continuation argument is an
elementary instance of that method. The result developed here is the
complete real zero classification for this specific odd-denominator
harmonic interpolation, together with quantitative zero asymptotics.
No exhaustive literature priority claim is made.

\section{Entire continuation and a zero-free real ray}

Define, for $t\ge0$,
\begin{equation}\label{sorder:def:K}
 K(t)=\frac{e^{-t}\log(1+e^{-2t})}{1+e^{-2t}}
     =\frac{\log(2\cosh t)-t}{2\cosh t}.
\end{equation}
The second equality holds on the real line; near zero both sides have
the same analytic germ. The generating function of the harmonic
numbers gives
\[
 -K(t)=\sum_{n\ge0}(-1)^nH_n e^{-(2n+1)t}\qquad(t>0).
\]
Consequently, initially by absolute convergence for $\Re s>1$,
\begin{equation}\label{sorder:eq:Mellin}
 S(s)=-\frac1{\Gamma(s)}\int_0^\infty t^{s-1}K(t)\,dt.
\end{equation}
The original series is locally uniformly convergent for $\Re s>0$:
summation by parts applies because $H_n(2n+1)^{-s}\to0$ and the
successive differences are $O_C(n^{-1-\sigma}\log n)$ on every compact
set $C\subset\{\Re s\ge\sigma>0\}$. Thus \eqref{sorder:eq:Mellin}
represents \eqref{sorder:def:S} throughout that half-plane.

\begin{proposition}\label{sorder:prop:entire}
The function $S$ extends to an entire function, real on the real axis.
For every integer $j\ge0$,
\begin{equation}\label{sorder:eq:negative-derivative}
 S(-j)=(-1)^{j+1}K^{(j)}(0).
\end{equation}
Moreover $S(s)<0$ for every real $s\ge-2$.
\end{proposition}
\begin{proof}
The function $K$ is analytic near zero and is $O(e^{-3t})$ at positive
infinity, as are each of its fixed-order derivatives. Subtract its
Taylor polynomial on $[0,1]$. The Mellin integral then continues
meromorphically across any finite part of the plane, with possible
simple poles at $s=-j$ and residues $K^{(j)}(0)/j!$.
The reciprocal gamma function has a simple zero there with leading
coefficient $(-1)^j j!$. This cancels every pole and gives
\eqref{sorder:eq:negative-derivative}. The construction uses real
Taylor coefficients and a real kernel, proving the reality assertion.

For the sign assertion we show $K''(t)>0$ on $[0,\infty)$. Set
$y=e^{-2t}\in(0,1]$. Direct differentiation gives
\begin{equation}\label{sorder:eq:Ksecond}
 K''(t)=\frac{\sqrt y}{(1+y)^3}
 \left((y^2-6y+1)\log(1+y)-4y^2+8y\right).
\end{equation}
If $y^2-6y+1\ge0$, the bracket is positive. Otherwise
$\log(1+y)<y$, so the bracket is strictly greater than
\[
 (y^2-6y+1)y-4y^2+8y=y(y-1)(y-9)\ge0.
\]
Two integrations by parts in \eqref{sorder:eq:Mellin} give
\begin{equation}\label{sorder:eq:twice-parts}
 S(s)=-\frac1{\Gamma(s+2)}
       \int_0^\infty t^{s+1}K''(t)\,dt,
       \qquad \Re s>-2.
\end{equation}
The integrations are first justified for $\Re s>0$; analyticity then
extends the identity to the stated half-plane. For real $s>-2$ all
factors in the integral and the gamma denominator are positive, so
$S(s)<0$. At the endpoint,
$S(-2)=-K''(0)=(\log2-1)/2<0$.
\end{proof}

\begin{corollary}[A complex zero-free half-plane]\label{sorder:cor:right-half-plane}
For every $s\in\mathbb C$ with $\Re s\ge3$,
\[
 |3^sS(s)+1|<\frac56.
\]
In particular $S$ has no zero in that half-plane.
\end{corollary}
\begin{proof}
The $n=1$ term of $3^sS(s)$ is $-1$, and absolute convergence gives
\[
 |3^sS(s)+1|\le27\sum_{n=2}^\infty\frac{H_n}{(2n+1)^3}.
\]
For $n\ge6$, use $H_n\le1+\log n$ and
$(2n+1)^3>8n^3$. The function $(1+\log x)x^{-3}$ decreases for
$x\ge5$, so the tail from $n=6$ is at most
$(3+2\log5)/800<7/800$. Therefore
\[
 \sum_{n=2}^\infty\frac{H_n}{(2n+1)^3}
 <\sum_{n=2}^5\frac{H_n}{(2n+1)^3}+\frac7{800}
 =\frac{122481517493}{3993750684000}.
\]
Multiplication by $27$ gives a rational number strictly below $5/6$.
\end{proof}

\subsection{Exact values at the negative integers}

Define the generalized secant numbers by
\[
 (\sec z)^\alpha=\sum_{m\ge0}E_{2m}(\alpha)\frac{z^{2m}}{(2m)!},
 \qquad E_{2m}=E_{2m}(1)>0.
\]
A prime on $E_{2m}(\alpha)$ will mean differentiation with respect
to $\alpha$. Comparison of the even and odd parts of \eqref{sorder:def:K}
gives the following exact formulas.
\begin{corollary}\label{sorder:cor:negative-values}
For $m\ge0$,
\begin{align}
 S(-2m)&=\frac{(-1)^m}{2}
       \left(E'_{2m}(1)-E_{2m}\log2\right),\label{sorder:eq:negative-even}\\
 S(-2m-1)&=\frac{(-1)^{m+1}}2(2m+1)E_{2m}.
 \label{sorder:eq:negative-odd}
\end{align}
For $m\ge1$, $(-1)^mS(-2m)>0$.
\end{corollary}
\begin{proof}
The even part of $K$ is
$\tfrac12\sech t\,[\log2+\log\cosh t]$ and its odd part is
$-\tfrac12t\sech t$. Use
$\partial_\alpha(\sech t)^\alpha|_{\alpha=1}
=-\sech t\log\cosh t$ and $\sech t=\sec(it)$ in
\eqref{sorder:eq:negative-derivative}. This proves both identities.
The Taylor coefficients of $\log\sec z$ in positive even degrees are
positive: its derivative is $\tan z$, whose positive coefficients
follow recursively from $u'=1+u^2$, $u(0)=0$.
Therefore, when $m\ge1$, $E_{2m}(\alpha)$ is a nonzero polynomial
in positive powers of $\alpha$ with nonnegative coefficients.
It follows that $E'_{2m}(1)\ge E_{2m}$, proving the sign because
$\log2<1$.
\end{proof}
For example,
\[
\begin{array}{c|rrrrrrr}
 s&0&-1&-2&-3&-4&-5&-6\\\hline
 S(s)&-\frac{\log2}{2}&-\frac12&\frac{\log2-1}{2}&\frac32&
 4-\frac52\log2&-\frac{25}2&\frac{61\log2-121}{2}
\end{array}
\]
The sign change between $-2m-1$ and $-2m$ already proves the existence
of at least one zero in each interval. It gives neither uniqueness nor
exclusion of other real zeros; those require the next argument.

\section{A Fourier--beta identity}

For real $\alpha>0$ define
\[
 B_\alpha(s)=\frac1{\Gamma(s)}\int_0^\infty
        \frac{t^{s-1}}{(2\cosh t)^\alpha}\,dt,
 \qquad\Re s>0.
\]
In particular $B_1(s)=\beta(s)$. From the kernels, or from the
binomial series and $(\alpha)_n$ differentiated at $\alpha=1$,
\begin{equation}\label{sorder:eq:parameter-derivative}
 S(s)=s\beta(s+1)+\left.\partial_\alpha B_\alpha(s)\right|_{\alpha=1}.
\end{equation}
This identity extends wherever its two sides are continued.

\begin{lemma}\label{sorder:lem:Fourier}
For $\Re s<1$,
\begin{equation}\label{sorder:eq:Fourier-B}
 B_\alpha(s)=\frac{\cos(\pi s/2)}{2\pi\Gamma(\alpha)}
 \int_0^\infty x^{-s}
 \Gamma\left(\frac{\alpha+ix}{2}\right)
 \Gamma\left(\frac{\alpha-ix}{2}\right)\,dx.
\end{equation}
On the left this means analytic continuation. Differentiation with
respect to $\alpha$ is valid locally uniformly near $\alpha=1$.
\end{lemma}
\begin{proof}
The substitution $u=e^{2t}$ and the beta integral show
\[
 \int_{-\infty}^\infty\frac{e^{ixt}}{(2\cosh t)^\alpha}\,dt
 =\frac{\Gamma((\alpha+ix)/2)\Gamma((\alpha-ix)/2)}{2\Gamma(\alpha)}.
\]
Fourier inversion and the Mellin transform of the cosine prove
\eqref{sorder:eq:Fourier-B} first in $0<\Re s<1$. For a direct
justification of the interchange, insert $e^{-\varepsilon t}$ before
the Mellin integral, integrate in $t$, and then send
$\varepsilon\downarrow0$. The resulting kernels are bounded near
zero by a constant times $x^{-\Re s}$, and the gamma product decays
exponentially at infinity. Dominated convergence applies.
The right side is holomorphic in $\Re s<1$, so it supplies the
claimed continuation. Differentiation in $\alpha$ only introduces
digamma factors growing logarithmically at infinity; the same
domination applies on a compact positive $\alpha$-interval.
\end{proof}

For real $q>0$ put
\begin{align}
 J(q)&=\int_0^\infty x^q\sech(cx)\,dx,
 \label{sorder:def:J}\\
 A(q)&=\int_0^\infty x^q\sech(cx)L(x)\,dx,
 \label{sorder:def:A}\\
 D(q)&=\int_0^\infty x^q\sech(cx)\tanh(cx)\,dx,
 \label{sorder:def:D}\\
 R(q)&=\frac{A(q)}{cD(q)}.
 \label{sorder:def:R}
\end{align}
The letter $D$ here denotes a positive integral, not differentiation.

\begin{proposition}\label{sorder:prop:reflected}
For $q>0$,
\begin{equation}\label{sorder:eq:reflected}
 S(-q)=\frac{cD(q)}2
       \left(R(q)\cos(cq)-\sin(cq)\right).
\end{equation}
Moreover,
\begin{equation}\label{sorder:eq:moments}
 J(q)=2\Gamma(q+1)c^{-q-1}\beta(q+1),\qquad
 D(q)=\frac q c J(q-1).
\end{equation}
\end{proposition}
\begin{proof}
At $\alpha=1$ the reflection formula gives
$\Gamma((1+ix)/2)\Gamma((1-ix)/2)=\pi\sech(cx)$.
The logarithmic $\alpha$ derivative of the full gamma quotient in
\eqref{sorder:eq:Fourier-B} is $L(x)$.
Thus the derivative term in \eqref{sorder:eq:parameter-derivative}
at $s=-q$ is $\tfrac12\cos(cq)A(q)$.
Applying the same Fourier formula to $\beta(1-q)$ gives
$\beta(1-q)=\tfrac12\sin(cq)J(q-1)$.
Since $D(q)=qJ(q-1)/c$ by integration by parts, the remaining term is
$-\tfrac12 c\sin(cq)D(q)$. This proves \eqref{sorder:eq:reflected}.
The expansion $\sech(cx)=2\sum_{n\ge0}(-1)^ne^{-(2n+1)cx}$
gives the stated formula for $J(q)$ for $q>0$ by absolute integration.
The integration-by-parts identity for $D$ is valid for every $q>0$.
\end{proof}

\section{The derivative bound that forces uniqueness}

The next estimate is the central analytic input. Its numerical constant
is deliberately coarse; it is small enough for the tangent comparison
and needs no numerical certification.

\begin{lemma}\label{sorder:lem:r}
Let $r(x)=L(x)\coth(cx)/c$ for $x>0$. Then
\begin{equation}\label{sorder:eq:r-derivative}
 0<xr'(x)\le a+\frac b x,\qquad
 a=\frac2c=\frac4\pi,\qquad
 b=\frac{2+2\log2}{c^2}.
\end{equation}
\end{lemma}
\begin{proof}
The digamma series gives
\begin{align}
 L(x)&=-2\log2+2\sum_{\substack{k\ge1\\k\text{ odd}}}
          \frac{x^2}{k(k^2+x^2)},\label{sorder:eq:L-series}\\
 L'(x)&=4x\sum_{\substack{k\ge1\\k\text{ odd}}}
          \frac{k}{(k^2+x^2)^2}>0.\label{sorder:eq:L-prime}
\end{align}
These series converge locally uniformly in $x$. Also
\[
 L'(x)=\int_0^\infty\frac{u\sin(xu)}{\sinh u}\,du.
\]
Writing $h(u)=u/\sinh u$, integration by parts yields
\[
 xL'(x)=1+\int_0^\infty h'(u)\cos(xu)\,du.
\]
Here $h$ decreases from $1$ to $0$, so
$\int_0^\infty|h'(u)|\,du=1$ and $0<xL'(x)\le2$.

Now
\[
 r'(x)=\frac{L'(x)}c\coth(cx)-L(x)\csch^2(cx).
\]
This is positive whenever $L(x)\le0$. If $L(x)>0$, comparison term
by term in \eqref{sorder:eq:L-series}--\eqref{sorder:eq:L-prime}
gives
\[
 L(x)<L(x)+2\log2\le\frac{1+x^2}{2}\,xL'(x).
\]
It follows that
\[
 r'(x)>\frac{L'(x)}c\coth(cx)
 \left(1-\frac{cx(1+x^2)}{\sinh(2cx)}\right)>0.
\]
The last inequality follows from
$\sinh(\pi x)\ge\pi x+(\pi x)^3/6>cx(1+x^2)$.
Finally $L(x)\ge-2\log2$, $\coth y\le1+1/y$ and $\sinh y\ge y$
give
\[
 xr'(x)\le\frac2c\left(1+\frac1{cx}\right)
       +\frac{2\log2}{c^2x},
\]
which is \eqref{sorder:eq:r-derivative}.
\end{proof}

For completeness we give the needed one-dimensional form of the
Brascamp--Lieb variance inequality \cite{BL}.
\begin{lemma}\label{sorder:lem:BL}
For a probability density proportional to $e^{-V(x)}$ on an interval,
with $V''(x)>0$,
\[
 \Var(f)\le\E\frac{f'(x)^2}{V''(x)}.
\]
Consequently
$|\Cov(f,g)|^2\le
\E(f'^2/V'')\E(g'^2/V'')$ whenever the displayed quantities exist.
\end{lemma}
\begin{proof}
It suffices first to work on a compact interval $[a,b]$ and subtract
the mean of $f$. Put
$u(x)=e^{V(x)}\int_a^x f(t)e^{-V(t)}\,dt$.
Then $u(a)=u(b)=0$ and $f=u'-V'u$. Integration by parts gives
\[
 \E f^2=-\E(f'u),\qquad
 \E f^2=\E u'^2+\E(V''u^2).
\]
Cauchy--Schwarz with weights $V''$ implies
$(\E f^2)^2\le\E(f'^2/V'')\E(V''u^2)
\le\E(f'^2/V'')\E f^2$. Division proves the variance bound.
Exhaust the original interval by compact intervals and use truncation
and convergence of the stated integrals. The covariance assertion
then follows from Cauchy--Schwarz and the two variance bounds.
\end{proof}

\begin{theorem}\label{sorder:thm:R}
For every real $q\ge2$,
\begin{equation}\label{sorder:eq:R-bound}
 R(q)>0,\qquad
 0<R'(q)<\frac1q\sqrt{\frac{181}{27}}<\frac\pi2.
\end{equation}
\end{theorem}
\begin{proof}
Give $(0,\infty)$ the probability density
\[
 d\nu_q(x)=\frac{x^q\sech(cx)\tanh(cx)}{D(q)}\,dx.
\]
Then $R(q)=\E_{\nu_q}r(X)$ and differentiation in $q$ gives
$R'(q)=\Cov_{\nu_q}(r(X),\log X)>0$, because both functions are
strictly increasing and the density has full support. Differentiation
is justified by absolute integrability of the additional logarithm.
Since $S(-2)=(\log2-1)/2$, \eqref{sorder:eq:reflected} gives
$A(2)=1-\log2>0$, whence $R(2)>0$ and $R(q)>0$ for all $q\ge2$.

The negative logarithm of the density, up to a constant, is
\[
 V_q(x)=-q\log x+\log\cosh(cx)-\log\tanh(cx).
\]
Its second derivative is
\[
 V_q''(x)=\frac q{x^2}+c^2\sech^2(cx)
     +4c^2\frac{\cosh(2cx)}{\sinh^2(2cx)}\ge\frac q{x^2}.
\]
Lemma~\ref{sorder:lem:BL}, applied to $r$ and $\log x$, therefore yields
\begin{equation}\label{sorder:eq:covariance-bound}
 R'(q)\le\frac1q\sqrt{\E_{\nu_q}(Xr'(X))^2}.
\end{equation}
There is no endpoint assumption hidden here: as $x\downarrow0$,
$r(x)=O(x^{-1})$, $xr'(x)=O(x^{-1})$, and the density is
$O(x^{q+1})$. All second moments and energy integrals used in the
lemma are finite for $q\ge2$. The tails are exponentially integrable.

For $j=1,2$, $\E_{\nu_q}X^{-j}$ decreases with $q$ by the same
covariance differentiation, now applied to a decreasing function and
$\log x$. At $q=2$ the two inverse moments coincide. Indeed,
\[
 D(0)=\frac1c,\quad D(1)=\frac1c,\quad D(2)=\frac{4G}{c^3},
\]
so
\[
 \E_{\nu_q}X^{-1},\ \E_{\nu_q}X^{-2}
 \le\mu:=\frac{\pi^2}{16G}<\frac34.
\]
Here $G>1-1/9=8/9$ by the alternating-series estimate and
$\pi^2<10$, so even $\mu<45/64$ suffices. The elementary inequalities
$\pi>3$ and $\log2<3/4$ give $a<4/3$ and $b<14/9$ in
Lemma~\ref{sorder:lem:r}. Consequently
\[
 \E_{\nu_q}(Xr'(X))^2
 \le a^2+2ab\mu+b^2\mu
 <\frac{16}{9}+\frac{28}{9}+\frac{49}{27}
 =\frac{181}{27}.
\]
This proves the first derivative bound. For $q\ge2$ its square is
at most $181/108<9/4<c^2$, proving the final strict inequality.
\end{proof}

\section{Complete classification of the real zeros}

\begin{theorem}\label{sorder:thm:all-real-zeros}
For each integer $m\ge1$, $S$ has exactly one real zero
\[
 \sigma_m\in(-2m-1,-2m).
\]
Every such zero is simple. These are all the real zeros of $S$.
The displacements $\varepsilon_m=\sigma_m+2m+1$ strictly decrease
with $m$. In particular, $\sigma_m-\sigma_{m+1}>2$.
\end{theorem}
\begin{proof}
There are no zeros on $[-2,\infty)$ by
Proposition~\ref{sorder:prop:entire}. Write $q=-s>2$.
On $(2m,2m+1)$, equation \eqref{sorder:eq:reflected} is equivalent to
\[
 \tan(c(q-2m))=R(q).
\]
The left side increases from $0$ to $+\infty$. The right side is
positive and finite, and the derivative of their difference is
$c\sec^2(c(q-2m))-R'(q)>0$ by Theorem~\ref{sorder:thm:R}.
There is therefore exactly one solution. At that solution the same
strict inequality and $\cos(cq)\ne0$ show that the derivative of
$S(-q)$ is nonzero, proving simplicity.

On $(2m+1,2m+2)$, $\cos(cq)$ and $-\sin(cq)$ have the same sign.
Since $R(q)>0$ and $D(q)>0$, \eqref{sorder:eq:reflected} cannot vanish.
At the integer endpoints one of the sine and cosine is zero and the
other is nonzero, while $R,D>0$. This exhausts $q>2$.

For the monotonicity write $q_m=2m+\delta_m$, where
$\delta_m\in(0,1)$ and $\tan(c\delta_m)=R(2m+\delta_m)$.
Increasing $m$ strictly increases the right side at every fixed
$\delta$. Since the difference between the left and right sides is
strictly increasing in $\delta$, its unique zero moves to the right.
Thus $\delta_{m+1}>\delta_m$, or
$\varepsilon_{m+1}<\varepsilon_m$, as claimed.
\end{proof}

\section{All-orders asymptotics of the zero displacement}

For $k\ge1$ define
\[
 d_k=\frac{(-1)^{k+1}2^{2k}B_{2k}(1/2)}{2k},
 \quad d_1=-\frac16,\quad d_2=-\frac7{60},\quad d_3=-\frac{31}{126}.
\]
Here $B_n(x)$ is the Bernoulli polynomial. Let
$q^{\underline j}=q(q-1)\cdots(q-j+1)$.

\begin{lemma}\label{sorder:lem:R-asymptotic}
For every fixed integer $K\ge1$, as real $q\to+\infty$,
\begin{equation}\label{sorder:eq:R-asymptotic}
 cR(q)=\psi(q+1)+\gamma-\log\pi
       +\sum_{k=1}^{K-1}\frac{d_kc^{2k}}{q^{\underline{2k}}}
       +O_K(q^{-2K}).
\end{equation}
In particular,
\begin{equation}\label{sorder:eq:R-first}
 cR(q)=\log(q/\pi)+\gamma+\frac1{2q}
                  -\frac{2+\pi^2}{24q^2}+O(q^{-3}).
\end{equation}
\end{lemma}
\begin{proof}
The ordinary sectorial digamma expansion \cite[\S5.11]{DLMF}, applied on the vertical
line with real part $1/2$, gives
\[
 L(x)=\gamma+\log(x/2)+\sum_{k=1}^{K-1}d_kx^{-2k}
                     +O_K(x^{-2K})\qquad(x\to\infty).
\]
The coefficient formula follows either from the shifted Bernoulli
expansion of $\psi(z+1/2)$ or by expanding the ordinary formula in
$1/(ix/2)$. The error may be bounded globally by $C_Kx^{-2K}$ for
all $x>0$: on $(0,1]$ the actual function $L$ is bounded, and each
subtracted power and logarithm is bounded by a constant times
$x^{-2K}$. Hence for $q>2K-1$,
\[
 A(q)=J'(q)+(\gamma-\log2)J(q)
       +\sum_{k=1}^{K-1}d_kJ(q-2k)+O_K(J(q-2K)).
\]
Use \eqref{sorder:eq:moments}. The first quotient, multiplied by $c$,
is
\[
 \frac{\beta(q+1)}{\beta(q)}
 \left(\psi(q+1)+\gamma-\log\pi+
                         \frac{\beta'(q+1)}{\beta(q+1)}\right).
\]
The $k$th remaining quotient is
$d_kc^{2k}\beta(q-2k+1)/(q^{\underline{2k}}\beta(q))$.
For fixed shifts, the beta ratios are $1+O_K(3^{-q})$, and the
logarithmic derivative is $O_K(3^{-q})$, directly from their
absolutely convergent Dirichlet series. The integrated error divided
by $qJ(q-1)/c$ is $O_K(q^{-2K})$. This proves
\eqref{sorder:eq:R-asymptotic}; expanding $\psi(q+1)$ gives
\eqref{sorder:eq:R-first}.
\end{proof}

\begin{theorem}\label{sorder:thm:zero-asymptotic}
Write
\[
 N=2m+1,\quad \sigma_m=-N+\varepsilon_m,\quad
 \Lambda=\log(N/\pi)+\gamma,\quad Q=\Lambda^2+c^2,
 \quad e_0=\frac1c\arctan\frac c\Lambda.
\]
Then $0<\varepsilon_m<1$ and
\begin{equation}\label{sorder:eq:zero-two-scales}
 \varepsilon_m=e_0-
       \frac{\tfrac12-e_0}{NQ}
       +O\left(\frac1{N^2\Lambda^2}\right)
       \qquad(m\to\infty).
\end{equation}
For every fixed $K\ge1$,
\begin{equation}\label{sorder:eq:zero-log-series}
 \varepsilon_m=
 \sum_{j=0}^{K-1}\frac{(-1)^jc^{2j}}{(2j+1)\Lambda^{2j+1}}
       +O_K(\Lambda^{-2K-1}).
\end{equation}
Thus the first three logarithmic terms are
$\Lambda^{-1}-\pi^2/(12\Lambda^3)+\pi^4/(80\Lambda^5)$.
The consecutive-zero gaps satisfy
\begin{equation}\label{sorder:eq:zero-gaps}
 \sigma_m-\sigma_{m+1}
 =2+\frac{2}{N(\Lambda^2+c^2)}
       +O\left(\frac1{N^2\Lambda^2}\right).
\end{equation}
\end{theorem}
\begin{proof}
The exact zero equation is
\begin{equation}\label{sorder:eq:epsilon-equation}
 c\cot(c\varepsilon_m)=cR(N-\varepsilon_m).
\end{equation}
Lemma~\ref{sorder:lem:R-asymptotic}, uniformly for $0\le\varepsilon\le1$,
gives
\[
 cR(N-\varepsilon)
 =\Lambda+\frac{\tfrac12-\varepsilon}{N}+O(N^{-2}).
\]
The solution of $c\cot(ce_0)=\Lambda$ is the stated $e_0$ for all
sufficiently large $N$. The right side of the preceding estimate
shows that the left side at $\varepsilon_m$ differs from $\Lambda$
by $O(N^{-1})$. Its inverse has derivative
$-1/(\Lambda^2+c^2)$ at $\Lambda$, so
$\varepsilon_m-e_0=O(N^{-1}\Lambda^{-2})$.
Taylor expansion of the left side about $e_0$, whose first derivative
is $-Q$ and second derivative is $2Q\Lambda$, now yields
\[
 -Q(\varepsilon_m-e_0)
 =\frac{\tfrac12-e_0}{N}+O(N^{-2}).
\]
The quadratic error is $O(N^{-2}/\Lambda)$ and is absorbed in the
displayed bound; the displacement inside the right side is even
smaller. Division by $Q$ proves \eqref{sorder:eq:zero-two-scales}.
Finally expand the arctangent. Its Taylor remainder gives the error
in \eqref{sorder:eq:zero-log-series}, while the inverse-index correction
is smaller than every fixed power of $1/\Lambda$.
For \eqref{sorder:eq:zero-gaps}, regard the displayed approximation
as a smooth function of real $N$. Its leading term has derivative
$e_0'(N)=-1/(N(\Lambda^2+c^2))$, and the derivative of the explicit
inverse-index correction is $O(N^{-2}\Lambda^{-2})$.
Take the difference at $N$ and $N+2$ in
\eqref{sorder:eq:zero-two-scales}; its two remainder values have the
same $O(N^{-2}\Lambda^{-2})$ bound.
\end{proof}

\paragraph{An algorithm for every inverse-index correction.}
Equation \eqref{sorder:eq:R-asymptotic} supplies a full expansion in
$1/q$. Substituting
$\varepsilon=e_0+\sum_{j\ge1}a_j(\Lambda)N^{-j}$ into
\eqref{sorder:eq:epsilon-equation} and comparing coefficients determines
$a_j$ recursively: the coefficient of the new unknown is always
$-Q\ne0$. The functions $a_j$ belong to
$\mathbb Q[c^2,\Lambda,e_0,Q^{-1}]$, with Bernoulli numbers as rational
coefficients, and $a_j=O_j(\Lambda^{-2})$.
For each fixed $J\ge1$ the truncation through $a_{J-1}N^{1-J}$ has
error $O_J(N^{-J}\Lambda^{-2})$. To justify this statement, first use
the full remainder in \eqref{sorder:eq:R-asymptotic} to bound the
residual after any fixed formal substitution by $O_J(N^{-J})$.
Derivatives of $c\cot(c\varepsilon)$ at $e_0$ are polynomials in
$\Lambda,c^2$ of degree at most the derivative order plus one.
Induction then gives $a_j=O_j(\Lambda^{-2})$ and controls all Taylor
remainders. The derivative of the exact implicit equation has magnitude
comparable to $Q$, by \eqref{sorder:eq:R-bound} and
$\varepsilon\asymp1/\Lambda$. The mean value theorem divides the
residual by $Q$, proving the stated error. This is an asymptotic
recursion; no convergence of the resulting infinite formal series is asserted.

